%%% Hlavní soubor. Zde se definují základní parametry a odkazuje se na ostatní části. %%%

%% Verze pro jednostranný tisk:
% Okraje: levý 40mm, pravý 25mm, horní a dolní 25mm
% (ale pozor, LaTeX si sám přidává 1in)
\documentclass[a4paper,oneside,12p]{report}
%\documentclass[a4paper,oneside,12p,twoside]{report}
\setlength\textwidth{155mm}
\setlength\textheight{237mm}
%\setlength\oddsidemargin{00mm}
\setlength\oddsidemargin{05mm}
\setlength\evensidemargin{05mm}
\let\openright=\clearpage


%% Vytváříme PDF/A-2u
\usepackage[a-2u]{pdfx}

%% Přepneme na českou sazbu a fonty Latin Modern
\usepackage[czech]{babel}
\usepackage{lmodern}
\usepackage[T1]{fontenc}
\usepackage{textcomp}

%% Použité kódování znaků: obvykle latin2, cp1250 nebo utf8:
\usepackage[utf8]{inputenc}

%%% Další užitečné balíčky (jsou součástí běžných distribucí LaTeXu)
\usepackage{amsmath}        % rozšíření pro sazbu matematiky
\usepackage{amsfonts}       % matematické fonty
\usepackage{amsthm}         % sazba vět, definic apod.
\usepackage{bbding}         % balíček s nejrůznějšími symboly
			   										% (čtverečky, hvězdičky, tužtičky, nůžtičky, ...)
\usepackage{bm}             % tučné symboly (příkaz \bm)
\usepackage{graphicx}       % vkládání obrázků
\usepackage{fancyhdr}				% možnost slylizovat záhlaví
\usepackage{fancyvrb}       % vylepšené prostředí pro strojové písmo
\usepackage{indentfirst}    % zavede odsazení 1. odstavce kapitoly
\usepackage[nottoc]{tocbibind} % zajistí přidání seznamu literatury,
\usepackage{icomma}         % inteligetní čárka v matematickém módu
\usepackage{dcolumn}        % lepší zarovnání sloupců v tabulkách
\usepackage{booktabs}       % lepší vodorovné linky v tabulkách
\usepackage{paralist}       % lepší enumerate a itemize
\usepackage{caption}				%	popisky
\usepackage{dirtree}				% strom souborů
\usepackage[bottom]{footmisc}  % poznámky pod čarou vespod
\usepackage{bibentry}
\nobibliography*

\usepackage{color}
\usepackage{forest}					% na rodokmeny
\usepackage{natbib}
%\setbibentrystyle{author}

\definecolor{pblue}{rgb}{0.13,0.13,1}
\definecolor{pgreen}{rgb}{0,0.5,0}
\definecolor{pred}{rgb}{0.9,0,0}
\definecolor{pgrey}{rgb}{0.46,0.45,0.48}

%%% Údaje o práci

\def\NazevSkoly{Gymnázium, Praha 6, Arabská 14}
% Název oboru včetně počátečního 'Obor'.
\def\NazevOboru{Dějepis}

% Název práce v jazyce práce (přesně podle zadání)
\def\NazevPrace{Židovké město v Třebíči}

% Název práce v angličtině
\def\NazevPraceEN{Jewish town in Třebíč}

% Název práce v němčině
\def\NazevPraceDE{Jüdische Stadt in Třebíč}

% Jméno auto
\def\AutorPrace{Havránek Kryštof 2.E}

% Rok odevzdání
\def\RokOdevzdani{2020}
% Měsíc odevzdání
\def\MesicOdevzdani{Květen}

% Vedoucí práce: Jméno a příjmení s~tituly
\def\Vedouci{PhDr. Lenka Dvořáková }

% Nepovinné poděkování (vedoucímu práce, konzultantovi, tomu, kdo
% zapůjčil software, literaturu apod.)
\def\Podekovani{%
\textbf{Poděkování}

TODO

PhDr. Lenka Dvořáková

Milan Krčmář (zjisit si titul!)
}

% Abstrakt (doporučený rozsah cca 80-200 slov; nejedná se o zadání práce)
\def\Abstrakt{%
TODO
}
\def\AbstraktEN{%
TODO
}
\def\AbstraktDE{%
TODO
}

% 3 až 5 klíčových slov (doporučeno), každé uzavřeno ve složených závorkách
\def\KlicovaSlova{%
{klíčová} {slova}
}
\def\KlicovaSlovaEN{%
{key} {words}
}


%% Balíček hyperref, kterým jdou vyrábět klikací odkazy v PDF,
%% ale hlavně ho používáme k uložení metadat do PDF (včetně obsahu).
%% Většinu nastavítek přednastaví balíček pdfx.
\hypersetup{unicode}
\hypersetup{breaklinks=true}

%% Definice různých užitečných maker (viz popis uvnitř souboru)
%%% Tento soubor obsahuje definice různých užitečných maker a prostředí %%%
%%% Další makra připisujte sem, ať nepřekáží v ostatních souborech.     %%%

%%% Drobné úpravy stylu

% Tato makra přesvědčují mírně ošklivým trikem LaTeX, aby hlavičky kapitol
% sázel příčetněji a nevynechával nad nimi spoustu místa. Směle ignorujte.
\makeatletter
\def\@makechapterhead#1{
  {\parindent \z@ \raggedright \normalfont
   \Huge\bfseries \thechapter. #1
   \par\nobreak
   \vskip 20\p@
}}
\def\@makeschapterhead#1{
  {\parindent \z@ \raggedright \normalfont
   \Huge\bfseries #1
   \par\nobreak
   \vskip 20\p@
}}
\makeatother

% Toto makro definuje kapitolu, která není očíslovaná, ale je uvedena v obsahu.
\def\chapwithtoc#1{
\chapter*{#1}
\addcontentsline{toc}{chapter}{#1}
}

% Trochu volnější nastavení dělení slov, než je default.
\lefthyphenmin=2
\righthyphenmin=2

% Zapne černé "slimáky" na koncích řádků, které přetekly, abychom si
% jich lépe všimli.
\overfullrule=1mm

%%% Makra pro definice, věty, tvrzení, příklady, ... (vyžaduje baliček amsthm)

\theoremstyle{plain}
\newtheorem{veta}{Věta}
\newtheorem{lemma}[veta]{Lemma}
\newtheorem{tvrz}[veta]{Tvrzení}

\theoremstyle{plain}
\newtheorem{definice}{Definice}

\theoremstyle{remark}
\newtheorem*{dusl}{Důsledek}
\newtheorem*{pozn}{Poznámka}
\newtheorem*{prikl}{Příklad}

%%% Prostředí pro důkazy

\newenvironment{dukaz}{
  \par\medskip\noindent
  \textit{Důkaz}.
}{
\newline
\rightline{$\square$}  % nebo \SquareCastShadowBottomRight z balíčku bbding
}

%%% Prostředí pro sazbu kódu, případně vstupu/výstupu počítačových
%%% programů. (Vyžaduje balíček fancyvrb -- fancy verbatim.)

\DefineVerbatimEnvironment{code}{Verbatim}{fontsize=\small, frame=single}

%%% Prostor reálných, resp. přirozených čísel
\newcommand{\R}{\mathbb{R}}
\newcommand{\N}{\mathbb{N}}

%%% Užitečné operátory pro statistiku a pravděpodobnost
\DeclareMathOperator{\pr}{\textsf{P}}
\DeclareMathOperator{\E}{\textsf{E}\,}
\DeclareMathOperator{\var}{\textrm{var}}
\DeclareMathOperator{\sd}{\textrm{sd}}

%%% Příkaz pro transpozici vektoru/matice
\newcommand{\T}[1]{#1^\top}

%%% Vychytávky pro matematiku
\newcommand{\goto}{\rightarrow}
\newcommand{\gotop}{\stackrel{P}{\longrightarrow}}
\newcommand{\maon}[1]{o(n^{#1})}
\newcommand{\abs}[1]{\left|{#1}\right|}
\newcommand{\dint}{\int_0^\tau\!\!\int_0^\tau}
\newcommand{\isqr}[1]{\frac{1}{\sqrt{#1}}}

%%% Vychytávky pro tabulky
\newcommand{\pulrad}[1]{\raisebox{1.5ex}[0pt]{#1}}
\newcommand{\mc}[1]{\multicolumn{1}{c}{#1}}


%% Titulní strana a různé povinné informační strany
\fancypagestyle{plain}{
\fancyhead[C]{}
\fancyhead[L]{Ročníková práce - Gymnázium, Praha 6, Arabská 14}
\fancyhead[R]{\textbf{Židovské město v Třebíči}}
\fancyfoot[L]{Vypracoval: Havránek Kryštof 2.E (Programování)}
\fancyfoot[C]{}
\fancyfoot[R]{\thepage}
\renewcommand{\headrulewidth}{0.4pt}
\renewcommand{\footrulewidth}{0.4pt}
}



\begin{document}



%%% Titulní strana práce

\pagestyle{empty}
\hypersetup{pageanchor=false}
%%\setlength\oddsidemargin{15mm}

\begin{center}

{\LARGE\bfseries\NazevSkoly}

\vspace{-18mm}
\vfill

{\LARGE\NazevOboru}

\vfill

\centerline{\mbox{\includegraphics[height=4cm]{../img/logo.png}}}

\vspace{-8mm}
\vfill

{\bf\Large ROČNÍKOVÁ PRÁCE}

\vfill


\vspace{15mm}

{\LARGE\bfseries\NazevPrace}


\vfill


Vypracoval: \hfill {\AutorPrace}

Vedoucí práce: \hfill {\Vedouci}

\vspace{15mm}
\MesicOdevzdani \ \RokOdevzdani

\end{center}



\newpage


\openright
%%\setlength\oddsidemargin{00mm}

\hypersetup{pageanchor=true}
\pagestyle{plain}
\pagenumbering{roman}
\vglue 0pt plus 1fill

\medskip\noindent
Prohlašuji, že jsem jediným autorem této práce, všechny citace jsou řádně označené a všechna použitá literatura a další zdroje jsou v práci uvedené.
Tímto dle zákona 121/2000 Sb. (tzv. Autorský zákon) ve znění pozdějších předpisů uděluji bezúplatně škole Gymnázium, Praha 6, Arabská 14 oprávnění k výkonu práva na rozmnožování díla (§ 13) a práva na sdělování díla veřejnosti (§ 18) na dobu časově neomezenou a bez omezení územního rozsahu.

\vspace{10mm}

\hbox{\hbox to 0.5\hsize{%
V ........ dne ............
\hss}\hbox to 0.5\hsize{%
Podpis autora
\hss}}

\vspace{20mm}

\newpage
%
% %%% Poděkování
%
\openright
%
\noindent
\Podekovani

\newpage


\openright

\vbox to 0.20\vsize{
\setlength\parindent{0mm}
\setlength\parskip{5mm}

Název práce:
\NazevPrace

Autor:
\AutorPrace

% Vedoucí práce:
% \Vedouci, \KatedraVedouciho

Abstrakt:
\Abstrakt

% Klíčová slova:
% \KlicovaSlova

\vss}\nobreak\vbox to 0.20\vsize{
\setlength\parindent{0mm}
\setlength\parskip{5mm}

% Opakování v angličtině.

Title:
\NazevPraceEN

Authors:
\AutorPrace

Abstract:
\AbstraktEN

\vss}\nobreak\vbox to 0.20\vsize{
\setlength\parindent{0mm}
\setlength\parskip{5mm}

% Opakování v němčině.

Titlel:
\NazevPraceDE

Autoren:
\AutorPrace

Abstrakt:
\AbstraktDE

\vss}

\newpage

\openright



\tableofcontents
\newpage


\chapter*{Úvod}
\addcontentsline{toc}{chapter}{Úvod}

TODO


\pagenumbering{arabic}
\setcounter{page}{1}

\chapter{Počátky města Třebíče}

\section{Charakteristika města Třebíče}

Dnes je Třebíč nepříliš velké město v jihovýchodní části kraje Vysočina.
I přes jeho velikost, však hrálo důležitou roli v historii okolí.
O tom svědčí i fakt, že má dokonce dvě památky zapsané v seznamu UNESCO a to Třebíčskou židovskou čtvrť a baziliku svatého Prokopa.
Počátky jeho historie jsou však, spíše než s národem židovským, spjaty s benediktýnským klášterem, dnešním zámkem v Třebíči.
Je možné, že na území Třebečského lesa se nacházelo i prehistorické osídlení, o tom však nemáme důkazy.

\section{Vznik kláštera}

Dějiny kláštera se začínají psát koncem 11. století.
Jednalo se o dobu kdy na českém i moravském území docházelo k řadě konfliktů mezi jednotlivými Přemyslovci a moravskými knížeti, právě ty měly za následek založení kláštera.
Nejvíce k založení kláštera napomohlo chování Břetislava II., který jako svého následníka jmenoval svého nevlastního syna Bořivoje II., místo brněnského knížete Oldřicha\footnote{ten měl být králem dle rodové tradice Přemyslovců}.
Kvůli této a dalším neshodám uspořádal roku 1099 král Břetislav tažení na brněnský hrad, což donutilo Oldřicha a jeho bratra Litolda, jenž spravoval znojemský úděl, k útěku.
Po dvou letech vlády Bořivoje se však Oldřich a Litold opět vrátili na své úděly, využili šance kdy se po smrti Břetislava II. Bořivoj II. musel vrátit do Prahy.

Již při cestě zpět si uvědomily, že západní hranice jejich údělu jsou příliš daleko od hradů v Brně a Znojmě.
Z tohoto důvodu v roce 1101 založili v prostorách Třebečského lesa téměř na hranicích s Přibyslavickem benediktýnský klášter.
Klášter měl sloužit jako předsunuté opevnění a kolonizační centrum.
Již v roce 1101\footnote{c. d. \bibentry{Uhlir1978}} byl postaven kostelík sv. Benedikta, který o tři roky později vysvětil pražský biskup Heřman.
Dále roku 1109\footnote{c. d. \bibentry{Uhlir1978}: 20} vysvětil biskup olomouckým Jan II. samotný klášterní chrám Nanebevztí Pany Marie.
Klášter ihned po založení dostával různá územní darování, a to jak od Oldřicha a Litolda, tak i od jejich nástupců.
Samotný Litold (\CrossOpenShadow 1112) a Oldřich (\CrossOpenShadow 1113) jsou dle slov Tomáše Pěšinky\footnote{Tomáš Petřinka z Čechorodu -- biskup, dějepisec, autor díla Mars Moravicus} pochováni v klášteře.

\begin{figure}[h]
	\centering
	\begin{forest}
		for tree={
    	child anchor=west,
    	parent anchor=east,
    	grow=east,
    	draw,
    	anchor=west,
    	edge path={
     		\noexpand\path[\forestoption{edge}]
       		(.child anchor) -| +(-5pt,0) -- +(-5pt,0) |-
       		(!u.parent anchor)\forestoption{edge label};
    		},
  	}
  	[Břetislav I.
    	[Vrastislav II.
				[Břetislav II.
						[Břetislav \textbf{D\footnotemark}
						]
						[Bořivoj II. (nevlastní syn)
						]
					]
				]
			[Kondrád I. Brněnský
				[Oldřich Brněnský \textbf{D}
					[Vrastislav \textbf{D}
					]
					[Vladislav
					]
				]
				[Litold Znojemský \textbf{D}
					[Kondrád II. Znojemksý \textbf{D}
					]
				]
			]
		]
		\end{forest}
		\caption[Rodokmen panovníků, vlastní tvorba]{
		\centering
			Rodokmen panovníků figurujících v rané historii kláštera, písmeno \textbf{D} značí knížata, která darovala území klášteru
		}
\end{figure}
\footnotetext{Dle opisu Kroniky Kosmovy, daruje vsi Břetislav, nejspíše se jedná o tohoto}

Bohužel je jen málo pramenů z této doby, jako jediný zdroj o rané historii kláštera máme Třebíčský (Brněnský) opis Kosmovy Kroniky\footnote{c. d. \bibentry{Uhlir1978}: 18}.
Do něho byla klášterním mnichem připsaná zmínka k roku 1115 o založení kláštera.
Mnich také uvádí jaké vsi byly postupně darovány klášteru.
Údaje jsou však sporné a pravděpodobně vycházel z pozdějších zdrojů, neboť počet darovaných vsí je dle vpisku sedmdesát tři, což značně převyšuje i pozdější čísla.
Ve vpisku jsou i uvedeny vsi u kterých je pozdější připojení písemně doloženo.
Existovala také klášterní kronika z roku 1201, jejímž autorem byl opat Tiburcius, ta se však ztratila.

\section{Tržní osady okolo kláštera}

\subsection{Podklášteří}

Vhodná poloha kláštera (procházelo okolím kláštera řada cest) a existence hradeb přispěly na přelomu 12. a 13. století k vzniku tržní osady.
Ta se samovolně začínala formovat přímo pod klášterem na levém břehu řeky, klášterní hradby nebyly totiž dost rozsáhle aby celou osadu pojaly.
Nově vzniklá obec Podklášteří ihned začala přitahovat kupce, ale i židy z širokého okolí, což vedlo k jejímu rychlému růstu.
Za nedlouho tak začala být omezována svým umístěním, byla totiž obehnána skálami a ohrožována záplavami řeky Jihlavy.
Podklášteří proto postupně ztrácelo na svém významu jakožto tržní osada a mnozí kupci odešli, zůstaly zde však židé, kteří položily základ židovskému městu\footnote{c. d. \bibentry{Uhlir1978}: 22}.

\subsection{Vznik Starého města}
%% přepsat
V 1. třetině 12. století vznikla na příkaz klášterní vrchnosti nová tržní osada, dnešní Staré město.
Ta se tentokrát nacházela na pravém břehu řeky Jihlavy, což umožňovalo více prostoru pro rozvoj.
Ihned po vytvoření se také začala tvořit aglomerace okolo nové osady, můžeme zmínit ves na Polánce, či vsi Kosovice, Arklebice a Horku, jenž opat Lukáš roku 1225 vyměnil se znojemskou šlechtičnou za ves Oslavany, které se nacházely dál od kláštera.
Stejně jako Podklášteří byla i nová osada omezena svým umístěním, což s postupem času opět začalo brzdit vývoj na území.
Proto ještě v 1. polovině 13. století bylo založeno Nové město.


\section{Založení města Třebíč}
Město Třebíč bylo podobně jako Staré město založeno z vůle opata, jenž si vyžádal povolení u krále Přemysla Otakara II.
Z toho důvodu se také o městu Třebíč hovoříme v této době jakožto o poddanském městu náležící církevní vrchnosti.
Opat dále musel vést vyšší soudy, či rozhodoval o manství\footnote{nápravníci, lidé disponující svobodným lénem}.
Opět se jednalo o strategický krok, města totiž disponovala různými privilegii, jako právo trhu či právo hradební, právem vářečným již disponovala Staržečska\footnote{Jiné označení Starého města, dnes čtvrť v Třebíči}.
Prví písemná zmínka o městu Třebíč pochází z roku 1277, v níž opat Martin oceňuje lokátora Heřmana\footnote{c. d. \bibentry{Uhlir1978}: 27}, který město Třebíč založil.

\subsubsection{Přestavba kláštera}

Během 13. století dochází k rozsáhle přestavbě kláštera.
Byly vztyčeny definitivní klášterní budovy a roku 1260 byla dokončena stavba klášterního chrámu.
Společně s přestavbou budov bylo přestavěn a vylepšen i hradební systém okolo kláštera.
Bohužel během Bitvy o Třebíč\footnote{12. květen — 9. červen 1468, součást česko-uherských válek (Druhá válka husitská)} byla většina staveb zničena či poškozena.
Postupem času vznikaly i další budovy patřící klášteru rozeseté po celém klášterem spravovaném území.

\chapter{Pozdější vývoj města Třebíče}

Město Třebíč se témeř okamžitě dostalo do problémů, díky bitvě na Moravském poli (26.srpen 1278) zavládla na Moravě nestabilní vláda a rostla moc menších knížat.
Tito knížata poté pořádala loupeživé výpravy se snahou zmocnit se královského majetku, s tím že ohroženo byla i samotný Třebíč. % 38
Situaci zachránil až nový král Jan Lucemburský, který sám Třebíč dvakrát navštívil.
Město poté dále prosperovalo a jeho význam pomalu přerostl i samotný klášter.
Výrazně pozdějšímu vývoji pomohlo privilegium markraběte Karla (Karel IV), díky kterému město získalo pevný právní řád.

Do dalších problémů se dostává město po smrti Jana Jindřicha (\CrossOpenShadow 1378) díky konfliktům Joště a Prokopa, kteří usilovali o kontrolu nad Moravou.
Tyto konflikty vyústily až v jakousi domácí válku, což vyvrcholilo koncem 14. století v těžkou hospodářskou krizi.
Situaci uklidnil až král Václav IV. roku 1413, doba klidu však neměla dlouhého trvání, protože již roku 1415 se začal stupňovat boj proti církvi, díky kázání mistra Jana Husa. % 43

\section{Třebíč a hustiská okupace}

O samotném vývoji husitství v Třebíči nemáme žádné prameny, je však známo, že Třebíč udržovala vztahy se moravskými šlechtici, jež byli myšlence nakloněni.
Během Žižkova tažení na Moravu se však třebíčský klášter dostal pod kontrolu husitů, z kláštera poté podnikaly výboje dál do okolí.
Roku 1430, však Prokop Veliký město dobývá zpět, konečně se situace uklidnila až v roce 1435.
Během hustiských válek došlo k výraznému hospodářskému úpadku, ze kterého se klášter dostával celá desetiletí. % 49-54

\section{Třebíč během Česko-uherských válek}

Třebíč měl po husitských válkách zájem o mír a byl tedy věrný králi Jiřímu z Poděbrad.
Již 12. května 1468 se Třebíč dostala do obležení, a o dva dny později byla Matyášovými vojsky dobyta.
Další boje se zaměřily na samotný klášter, ten obývalo asi čtyři tisíce lidí.
6. června se česká vojska pokusila klášter zachránit, ale bez úspěchu, 15. června se posádka v klášteře definitivně vzdala.
Po celých sedm let se bylo město opuštěné. % 54 - 56

\section{Vláda Perštejnů}

Na obnovách se začalo pracovat již roku 1474, s tím že nové město začalo vzrůstat na troskách starého.
Celková restaurace včetně klášterního velkostatku však nebyla možná díky obrovským dluhům, jenž klášter čelil.
Nejvíce pomohl Vilém z Perštejna\footnote{bohatý rod, ve službách Jagellovců, třebíčský statek je pod jeho kontrolou}, ten zavedl rybnikářství a započal opravy samotného kláštera.
Roku 1507 Vilém dělí své území mezi syny, ačkoliv apeluje ve své závěti na nutnost oprav různých klášterů v okoli o tom třebíčském nepadla ani zmíňka, existence benediktínů mu bránila v převedení třebíče do dědičného vlastnictví.
Déle se Třebíče ujímá Jan z Perštejna, který pokračuje v otcově politice, a na krátkou dobu město přesouvá pod nadvládu Jana Jetřicha, který jako první světský vládce razantně zasahuje do vedení města.
Samotný Jan z Perštena přiděluje městu další práva, z nichž je nejdůležitější právo mílové, které značně zlepší ekonomické postavení města.
Obnoven byl na žádost městké rady i vypovídací dekret Jana Jetřicha Černohorského, s cílem omezit konkurence schopnost židovské populace ve městě.
Po Janu(\CrossOpenShadow 1548) se vlády ujímá syn Vratislav, ten si již Třebíč převádí do dědičného vlastnictví. % 61 - 68/71
% režijní feudální velkostatek

\section{Vláda Osovců}
Mezi polovinou 16. století a rokem 1613 je Třebíč državou Osovců, kdy Burian Osovský odkupuje Třebíčské panství.
Vzhledem k ceně třebíčského panství se však Burian zadlužil, dluh se snažil vykompenzovat produkcí piva a souběžně omezoval práva města -- právo na pronájem, ale i právo mílové.
Město proto podalo stížnost na zemský soud, na neuspěla, a tak se Třebíč odvolala k císaři.
Ferdinant I. poté 7. března 1560 obnovuje platnost všech privilégií, zároveň napomíná samotného Buriana, ten již dále nepodnikal tak razantní kroky proti městu. % 77-78

Roku 1568 se vlády nad městem ujímá Smil Osovský z Doubravice\footnote{Jeden z důležitých pramenů doby, své paměti si zapisoval do Codexu Dubraviciana}
Smil musel čelit dluhům jeho otce, ty částečně vyřešil sňatky, ale nejvíce mu pomohla jeho hospodářská prozíravost.
Ta značně zlepšila hospodářské postavení města, vše však bylo na úkor poddaných a to v podobně vyšších rent, či delších robotních povinností.
Vykazoval však snahu postarat se o sirotky, zlepšit zabezpečení před požáry, lépe zabezpečit rektory, podpora městského špitálu, či mírnější tresty -- například za pytlačení nejsou vydloubnuta oka.
Roku 1573 dále vydal Zřízení selské, kde nabádá k poslušnosti vůči úředníkům a zakazuje sběr dříví a pastvu v panských lesích.
Roku 1582 mělo město právo pečetit červeným voskem a roku 1583 bylo vydáno městské zřízení -- město je zpravováno pánovými instrukcemi, neb měšťané ztratil zájem o správu. % přepsat
Je zajímavé, že touto dobou přicházejí do Třebíče i měšťané slovenské národnosti.

Třebíč tedy byla počátkem 17. století opět prosperujícím městem s rozvinutým řemeslem a jeho vládce již nebyl zadlužen rozvinutým řemeslem a jeho vládce již nebyl zadlužen, během vlády Osovců také došlo k přestavbě kláštera na zámek, jímž je i dnes. %78-96

Po úmrtí Smila se vlády ujímá po něm zbylá vdova Kateřina z Valdštejna, ta se sice provdala za Karla staršího ze Žertonína, ale Třebíč si nechala pod svojí kontrolou.
O politické příslušnosti Kateřiny nejsou sice písemné zmínky, víme však že k Třebíči se chovala mírně a udržovala dobré vztahy s měšťany. %  101-102

% purktretn9 knihy
\subsection{Židovské osídlení}

V urbářích\footnote{soupis pro evidenci platů a dávek} se dočteme i o stavu židovské populaci ve městě.
Urubář z roku 1573, tedy počátku Smilovy vlády, hovoří o dvou židech v Stařečce a šesti v Podkláštří.
Urubář z roku 1629 o jedenácti židech v podklášteří.
Během 16 století ještě židé sousedily v podklášteří s běžnými měšťany.
K růstu docházel až během konce 16. století a hlavně začátku 17 století, kdy začali skupovat domy měšťanů v podklášteří.
Samotná ekonomická situace židovské populace byla srovnatelná se zbytkem měšťanů, najdeme i zmínky i zadluženosti některých židů.
V urbářích se neobjevují ani jako věřitelé, až na jednu zmínku, či nejsou zaznamenány žádné ilegalní obchody.
Můžeme se domýšlet že již během 16. století existovala v Třebíči synagoga, protože počátkem 17. století vzniká synagoga s přízviskem nová. % 89

\section{Třicetiletá válka}

I přes politický nátlak od krajan; zůstal Žertonín nestranný a defenestraci Čechů neschvaloval jako adekvátní.
Žertonín ani nepodpořil nového krále Fridricha V. (zimní král), v jeho úspěch totiž nevěřil, naopak samotná Třebíč mu nabídla nocleh a předala dary při jeho cestě do Brna.
O dalším bředbělohorským vývojem nejsou žádné prameny, víme však, že Třebíče se žádná z událostí války příliš netýkala a  nikdy nebyla ohrožena.
Karel starší ze Třebíče také zaslíbil ochranu Třebíče, ta vyvrcholi 9. 12. 1620 kdy se sám císař Ferdinant II. zaručil, že císařská armáda město nepoškodí a bude ho chránit.
Mezi 14.-15. prosincem 1620 tak městem jen prošla císařská armáda a zanechala zde ochranou posádku.
Roku 1621 proto Třebíč dokázala vyhovět prosbám o pomoc Jihlavy a Znojma, která pod ochranou císaře nebyla. % 103-106

Počátkem ledna 1622 do města přijíždí Marradasův jezdecký pluk císařské armády o síle pěti set jezdců, který představuje pro Třebíč výraznou zátěž. %
Pluk se také nijak nehrnoul na bojiště a v městě pobyl až do října.
Do nebezpečí se město dostalo poté co Gábor Bethlen začal podnikat nájezdy na moravské území a pustošit místní města.
Třebíč prosila o pomoc, nebezpečí však bylo zažehnáno, když Gábor začal mírová jednání.
Mezitím do města však dorazily dva regimenty slezského vojska o síle tří tisíc vojáky, ti v Třebíči tábořily dvacet tři dní.
Roku 1626 došlo k dalšímu uherskému výboji, což vyvrcholilo opět v pobyt císařských vojsk v Třebíči. % 106 - 107

Další ránu městu zasadilo Obnovené zřízení zemské, poté už náboženskou svobodu, kterou Třebíč dříve disponovala, už nemohla šlechta nijak garantovat.
Brzy po vydání emigračního patentu (9. březen 1928) se proto Žertotín rozhodl pro emigraci, konkrétně do Vratislavy.
Z Vratislavy mohl stále spravovat území, které si ponechal, a Kateřina stále byla schopna dohlížet nad Třebíčí.
Novým držitelem Třebíče se na jaře roku 1928 stává Kateřin bratr Adam, samotná Kateřina roku 1638 zemřela a o dva roky dříve zemřel Karel starší. %
Jejich těla jsou pochována v Třebíči. % 108 -110

Samotný Adam nejevil o správu města velký zájem a přenechal ji synu Rudolfovi z Valdštejna.
Ten již neměl takové postavení politické jako jeho otec a také příliš nezasahoval do náboženských poměrů ve městě.
Otázky náboženství nepovažoval za příliš důležité a spíše se zajímal o otázky ekonomické, které řešil vysokými poplatky. % 109 - 112

Samotný vývoj třicetileté války během přechodu z nizozemské fáze do fáze švédské příliš Třebíč neovlivnil.
Až tažení Linharta Torstensona Třebíč postihla -- roku 1644 se musela podílet na vydržování a vystrojení vojska bránící Olomouc.
V březnu 1945 však do města přišla švédská jednotka, ta si vynutila 6500 zl. výpalného a obnovila v nekatolické bohoslužby, musela také platit poplatek švédské posádce v Jihlavě. % 112 - 113

\section{Vývoj po třicetileté válce}

Moravská šlecht pod vedením Jana Rottala začala vydávat jedno nařízení proti nekatolíkům za druhým, stejně se choval i císař se svými pateny z roku 1651.
Roku 1951 byla i sestavená komise s cílem provádět inspekci církevního majetku na Moravě.
I pro Třebíč tak druhá polovina 17. století znamenala období rekatolizace, až do 18. století však tajné nekatolictví v městě vydrželo.
Došlo také k obnovení cechů, které postupně eliminovali konkurenci a začali brzdit vývoj v městě. % 125 - 126

Samotné město se dostalo pod správu Zdislavy ze Sezimova Ústí, ta zastupoval syna Adama Františka.
O její zprávě nemáme příliš písemností, výmé však, že odpustila dluh kupcově rodině a vyhověla stížnosti židů, aby nebyly vystavováni násilí nestihnou li robotu.
30. června 1655 se již vlády ujímá syn Adam, ten se choval podobně jako jeho otec a otázce náboženství tak nepřikládal velkou hodnotu.
Již však nebyl schopen odolávat náporu z venčí a zvláště žalobě polenského děkana Kryštofa Kazimíra Burešovského, ta byla podložena řadou důkazů obsahovala však i polopravdy jako je existence 21 sekt v Třebíči, celkově však byl napsán na velmi vysoké úrovní.
Právě tato žaloba vyvolala na úřadech velkou reakci, což vyvrcholilo vytvořením \uv{inkviziční komice} zorganizované 2. prosince 1656 novoříšským proboštem Engelbertem a moravskobudějovickým Rudolfem Jindřichem ze Schaumburka.
Ta působila ve městě od 3. do 9. února 1657, za přítomnosti rychtáře, faráře a i samotného žalobce. % 126 - 128

Komise zjistila, že například v Podklášteří žilo 13 nekatolických rodin ku 7 katolickým a podobně na tom byl zbytek Třebíče.
Nejvíce dominovaly utrakvisté, vyskytovaly se však i luteráni, husité a  bratří.
Závěr komise je byl relativně jasný -- Třebíč je útočištěm nekatolíků, zvláště zběhlých z čech.
Komise si také stěžovala na chování židů, ti pracovali v neděli, o svátcích a během půstu prodávaly maso.
Samotná městská rada se proti nařčení z útlaku katolíků ohradila.
Děkan Burešovský a farář Slovatius navíc přidaly další stížnosti vůči Třebíči -- závěrečná zpráva tak pro Třebíč nebyla nikterak dobrá.  % 128 - 129

Již 8. května 1657 přichází do města protireformační komise vedená hejtmanem Tannazollem.
Její prvotní snahy však byly neúčinné celkově od listopadu 1658 do července 1659 se podařilo obrátit na \uv{pravou} víru jen 109 lidí, obecně lidé neměli přílišnou motivaci, pokuty pro nekatolíky nikdo nevymáhal.
Přelomem 50. a 60. let se však objevuje nová vlna emigrace -- 20\% šlechty. % 130 -131

Počátkem 60. let navíc začalo hrozit nebezpečí se strany turků, kteří pořádaly vpády na Moravu.
Samotná Třebíč však nebyla ohrožena a roku 1664 porážkou Turků u svatogotthardského průsmyku ohrožení pominulo nadobro. % 132 - 133

Roku 1666 umírá samotný Adam a Třebíč připadá Františkovy Augustinovy.
Za jeho vlády se dostává rekatolizace do nové fáze, František zároveň usiluje o zřízení kapucínského kláštera, jeho žádosti však nebylo vyhověno.
Umírá záhy a k moci se dostává jeho bratr Karel Ferdinand.
Ten se na rozdíl od jeho předchůdců zapsal jakožto dobrotivý pán.
Během  jeho vlády do města přicházejí i zástupci kapucínů a rychle se začleňují do městské společnosti. % 133 - 135

\section{Spory mezi měšťany a vrchností}

Za vlády Josefa I. docházelo k řadě sporů mezi měšťany a vrchností, ty zakončil až reskript císaře z 15. prosince 1708.
V něm císař ponechal Třebíči soudní pravomoc, zároveň však apeloval, aby se Valdštejn na jurisdikci podílel, zároveň zaručil měšťanům právo na svobodný trh. %

Během let 1724 až 1726 však došlo k dalším sporům mezi městem a vrchností -- Janem Josefem z Valdštejna, ty vyvrcholily narovnání roku 1926, které posunulo město dál k emancipaci. %

K dalšímu sporu došlo roku 1728, kvůli termínu tělesné člověčenství (vyjádření závislosti na feudálovi) jenž v jedné listině použil, tento spor měšťané prohrály a opět se objevila idea, že vztah měšťana a pána je vždy nevolnický.
Spor se dostal před tribunál ještě jednou díky kupci Janu Bernardovi Matesovi, který odmítl přijmout nevolnictví, načež mu bylo odebráno právo obchodovat, což bylo v rozporu s právem na svobodný trh.
Vrchnost tak byla nucena ustoupit.
Po smrti Jana Josefa se stal držitelem města František Heřman Arnošt z Valdštejna, ten se výrazu nevolnictví plně zříká a nepoužijí ho ani jeho následovníci. % 33 - 35

I za vlády Marie Terezie se stále spory vyskytovaly, jednalo se například o konflikt z roku 1744 ohledně přijetí Gallanta z Valašska (vrchnost byla pro, měšťané proti), spor ohledně nekompetentního rychtáře z roku 1750, či spor o repatriaci masného krejcaru z roku 1762.
Je důležité zmínit, že spory neprovázely jen Františka Heřmana, ale i jeho nástupce Františka Josefa Jiřího z Valdštejna \footnote{držitel mezi léty 1748 a 1760} a také jeho nástupce Emanuela Filiberta z Valdštejna \footnote{držitel mezi léty 1760 a 1775}.
Je zajímavé zmínit jaký měla příjem vrchnost (k roku 1779) -- z obyčejného měšťanstva 1294 zl. 30 kr., a z podklášterského židovstva 1141 zl. 36 kr.. %

\section{Vláda Marie Terezie}

Za vlády Marie Terezie se postupně začalo opouštět od tradiční ekonomice feudální společnosti a jít vstříc kapitalistickému fungování.
V Třebíči se vývoj odehrával podobně a na důležitosti nabyly tlkaci, ševci a koželuži, oproti tomu upadlo pro Třebíč dlouhou dobu charakteristické soukenictví, to se však dostalo do problémů z důsledku nedostupnosti vlny.
Situace okolo soukenictví vyvrcholila až v soudní spor.
Ve městě také vznikla papírna.
Můžeme také zmínit, že obchodem se živila většina židovských obyvatel v roce 1973 to bylo 55, v roce 1799 již 92.
Samotná průmyslová revoluce přišla do Třebíče trochu opožděně. % 13 - 29

Třebíč se stala obětí válek za děditství rakouské, kdy 10. února 1742 do města vtrhli pruští husaři a muselo jim být zaplaceno 12 000 zl. výpalného a 6000 zl. navíc od židů.
Dále 12. února dorazila armáda  krále Bedřicha Velikého o síle 12 000 mužů, kterou město muselo ubytovat a 24. února ještě Sasi, jejichž škody se odhadují na 90 000 zl.
Devastující byly také neúrody mezi léty 1765 a 1769, dále 1771 a 1772, ty vyvrcholily až v selské povstání z roku 1775. % 38

1. ledna 1755 přišla Třebíč o právo hrdelní, tomu předcházela její snaha omezit počet institucí, jenž tento trest mohly vynést.
Samotná Marie Terezie se však do vnitřní organizace měst příliš neangažovala, to však změnil její nástupce -- Josef II.. %

\section{Vláda Josefa II.}

První Josefův zásah je prostřednictvím nového dekretu o soudnictví z 27. prosince 1786, kdy je soudnictví soustředěno do kriminální soudů, jako které fungují magistráty krajských měst, v našem případě Jihalva.
Také došlo k zrušení městských rad a jejich nahrazení magistráty, jejichž zastupitelé -- purkmistr a rada volilo měšťanstvo prostřednictvím měšťanského výboru, samotné kandidáty však registrovala vrchnost, což jí dávalo větší moc.
Spory s vrchností však neustaly ani po vytvoření magistrátu (1. volby se konaly v březnu roku 1787), vrchnost také kromě schvalování kandidátů rozhodovala o výdajích města.
Docházelo také k silné germanizaci, a to jak ve spisech, tak v samotné správě města. % 36 - 37

\section{Napoleonské války}

Dne 17. srpna 1776 byl zveřejněn patent vyzývající občany k vojenské službě, samotné město dalo na válečnou snahu značný peněžní příspěvek.
To i přesto, že na přelomu 18. a 19. století bylo obecní hospodářství ve velmi špatném stavu, magistr nepřikročil k reparticím ani k vybrání peněz na obnovu kostelní věže.
Mezi 19. a 26. listopadem 1805 prošel městem první sbor bod vedením maršála Bernadottema, ten uvítal farář Dvořecký a nabídl Bernadottemu nocleh na zámku, ten se však spokojil s domem radního.
Za nedlouho však ve městě vypukl nedostatek jídla a francouzští vojáci museli táhnou dál.
Dále se ještě muselo město postarat o dvě divize, které sem přišli po bitvě u Slavkova a  pobývaly v Třebíči mezi 9. a 29. prosincem.
Celkové výdaje dosáhly 300 000 zl., ztráta však nikdy nebyla vykompenzována.
Situaci ještě zhoršily špatné úrody v letech 1807 a 1808. %

Po uzavření příměří po bitvě u Znojma Třebíč protínala demarkační linie, což způsobylo městu další finanční ztráty o celkové hodnotě 100 000 zl.
Navíc do města byly svezeni raněních obou bojujících stran, to mělo za následek tyfovou epidemii. %

\section{Vývoj po napoleonských válkách}

Další devastující událostí pro Třebíč byl státní bankrot z roku 1811, který obral malé živnostníky i o to málo co jim zbylo, mezi mocnými se také rychle rozšířila korupce.
V Třebíči navíc byly osočeni židé, že o bankrotu věděli předem a začali si skupovat majetek.
Do ohně ještě přilili neúrody v letech 1815 a 1817 a vše završila morová epidemie mezi dobytkem roku 1819, po které přestala fungovat většina řemeslné výrovy ve městě. % 39 - 40

Město utrpělo další ránu díky proti robotní vzpouře, ke které došlo roku 1821.
Vzbouřenci totiž byly odvedeni na zámek kde je ztrestali holemi, městem poté začaly kolovat letáky vyzývající k ukončení trestů s hrozbou žhářství.
Hrozby se stali skutečností a dne 3. května shořelo 100 domů v židovské obci a dalších 50 v Podklášteří. % 40

V červnu roku 1822 město opět zachvátil požár, který měl více devastující účinky než bitva o Třebíč.
Během červnových požárů shořelo přes 290 domů a to včetně kostela, radnice, fary a jedné školy.
Třebíč se poté opět dostala do finančních problému, které vyřešil až syndik Gottsmann.

Třebíči však nebylo dovoleno dlouhému oddychu, v roce 1830 byla Třebíč zasažena povodní, v roce 1832 a opět v roce 1836 vypukla epidemie cholery a v roce 1847 město opět zasáhl požár.

\section{Židovské město // nedokončené}

Na přelomu 17. a 18. století byly odkoupeny zbylé křesťanské domky v Podklášteřží a začal se pro něj používat termín židovské město, který předtím popisoval jen jeho část.
Podkálšteří fungovalo jako samostatná obec a nespadalo ani do stejné jurisdikce jako Třebíč.
V čele samotného židovského města stála rada v jejímž čele byl židovský rychtář, jenž byl volen občany města.

// strany 44 - 45, jen židovské město, více informací, zahrnu do kapitoly o samotném židovském městě.


// strany 49 - 64, jen řemeslná výroba a soudní spory ohledně ní, něco málo o zavádění strojů, ale vcelku irelevantní k mé práci, nehledě na to že od Marie Terezie nebyla žádná změna co se týče řemeslných oborů

\section{Národní obrození}
Snaha o národní obrození se v Třebíči začala objevovat ve 20. letech 19. století.
Většina obyvatel Třebíče byla česká, a případní cizinci (manželky) byly do českého prostředí rychle asimilováni.
Správu nad městem však držely němečtí úředníci a tak tempo obrození nebylo nijak valné.
První náznak změny představoval rodák Josef Chmela, který se stal velmi podlým spisovatelem a obdivovatelem Josefa Jungmana, v Třebíči pobýval o svých prázdninách kdy v hospodě zorganizoval malý debatní kroužek.
Jeho krátké návštěvy však neměly velký vliv, rozhodující osobností se stal až doktor Jan Miloslav Haněl.
Ten do Třebíče přišel jako mladý lékař roku 1835 z Prahy, a o devět let později založil čtenářský spolek a angažoval se také v divadelní činnosti, kterou pořádal spolu s českými studenty. % 67 - 69

Samotné události roku 1848 zastihly vedení Třebíče nepřipravené, situaci stabilizovalo až vytvoření národní gardy, to ale také vedlo k rozdělení měšťanstva na stranu německou, v jejímž čele stál Ignác Glas, a na stranu českou s dr. J. M. Hanělem v čele.
Zatím však národní garda pracovala společně s Němci. % 69 - 70

1. českou manifestací se stalo svěcení praporu národní gardy uspořádané dne 28. května 1849, při něm se zpívaly české písně a provolávala různá hesla.
Garda však neměla dlouhého trvání a již koncem roku 1849 musela z nařízení ukončit svoji činnost. % 70 - 72

20. října 1860 byl vydán Říjnový diplom, ve kterém byla zaslíbena konstituční vláda, načech 21. února 1861 byla vydána samotná ústava.
Obě události měšťané přijali s nadšením, a záhy se začala angažovat v Třebíči česká národní strana a vytvořila spolek Měšťanská beseda s dr. J. M. Hanělem v čele.
Jednalo se o jeden z prvních takových spolků na území Moravy.
Německá Třebíč jako odpověď na český spolek utvořila organizace Gesangverein, Turnverein a Schützverein (původní střelecký spolek). % 76

Období 40. let se však také neslo v duchu sociální nejistoty drobných výrobců, ta vycholila vystoupenímí namířenými hlavně proti židovské čtvrti.
Během let 48-49 však došlo k dočasnému zrovnoprávnění, což vedlo i k přesun zámožných židů z židovského města na náměstí.
To však mělo za následek další bouře proti židovskému obyvatelstvu, k těm došlo mezi 1. dubnem a 7. květem 1850.
K dalším nepokojům došlo v roce 1860, společně úpadkem pro Třebíč charakteristického soukenictví.
Ve městě také došlo k nedostatku peněz, obec tedy byla nucena vydat dočasné bankovky, ty se však velmi lehce padělaly.
Celá řada měšťanů se dostala do finančních problémů a byla nucena prodat své domy.
To mělo za následek další proti židovské nepokoje (1863 a 1866). % 77 - 78

\section{Prusko-rakouské války}

Během prosko-rakouské války došlo k poslednímu vzrůstu soukenictví, ekonomické situaci měšťanstva to však nijak nepomohl, zakázky z armády doprovázely i rekvizice, které město přišli až na 250 000 zl.
Třebíč se také dočkala opět nepřátelské invaze, a to 13. července, kdy do města přitáhla pruská pěchota, dělostřelectvo a jezdectvo, o síle 35 000 mužů.
Tato posádka v Třebíči pobyla do 16. prosince.
Dále město hostilo záložní sbory a zdravotní kolony.
Okupováno bylo až do uzavření příměří, poslední jednotky však odešli až 8. září.
Pruské jednotky do města navíc zavlekli choleru.

\section{Další vývoj}

20. srpan 1867 Třebíč dostala svého českého sartostu -- Ignáce Přerovského, zároveň radními příslušníky se stali členi Měšťanské besedy jako J. M. Haněl.
Vznikla také Občanská záložna, která měla fungovat jakožto opora českého národního hnutí.
K založení Sokola však v Třebíči nedošlo.

Při dalších volbách roku 1870 Německá stara vůbec nekandidovala a do obecního výboru tak byly zvoleni jen čeští funkcionáři.
\chapter*{Závěr}
\addcontentsline{toc}{chapter}{Závěr}

Cíl práce se i přes problémy podařilo splnit.
Bohužel se mi díky epidemii CoVid--19 nepodařilo město navštívit a získat tak mé osobní fotografie.
Ty však tvoří jen malou část práce.
Během nočního přejezdu jsem ve městě zastavil a čtvrť alespoň prošel.

Jediný problém ke zpracování představovalo období dějin města od revolučního roku 1848 po současnost.
Bohužel v tomto časovém rozsahu nejsou ani publikace ohledně samotného Třebíče.
Byla sice plánována kniha Třebíč: Dějiny města III., k jejímu vytvoření snad nikdy nedošlo, či vzniklo jen pár výtisků.
Obecně nedostatek středoškolskému studentu přístupných pramenů komplikuje popsat dějiny židovské komunity v Třebíči podrobně.

I tak si myslím, že je výsledkem této práce obstojný souhrnný pohled na historii židovské obce ve Třebíči.
Sám jsem díky této práci získal hluboký náhled do dějin nejen samotného Židovského města, ale i celého Třebíče a Moravy.
Pevně doufám, že se tohoto tématu ujme schopný historik a celé dějiny zmapuje podrobně od začátku do konce.


%%% Seznam použité literatury (bibliografie)
%%%
%%% Pro vytváření bibliografie používáme bibTeX. Ten zpracovává
%%% citace v textu (např. makro \cite{...}) a vyhledává k nim literaturu
%%% v souboru literatura.bib.
%%%
%%% Příkaz \bibliographystyle určuje, jakým stylem budou citovány odkazy
%%% v textu. V závorce je název zvoleného souboru .bst. Styly plainnat
%%% a unsrt jsou standardní součástí latexových distribucí. Styl czplainnat
%%% je dodáván s touto šablonou a bibTeX ho hledá v aktuálním adresáři.

\bibliographystyle{czechiso}    %% Autor (rok) s českými spojkami
% \bibliographystyle{plainnat}    %% Autor (rok) s anglickými spojkami
%\bibliographystyle{unsrt}       %% [číslo]

\renewcommand{\bibname}{Seznam použité literatury}

%%% Vytvoření seznamu literatury. Pozor, pokud jste necitovali ani jednu
%%% položku, seznam se automaticky vynechá.

\bibliography{literatura}

%%% Kdybyste chtěli bibliografii vytvářet ručně (bez bibTeXu), lze to udělat
%%% následovně. V takovém případě se řiďte normou ISO 690 a zvyklostmi v oboru.

% \begin{thebibliography}{99}
%
% \bibitem{lamport94}
%   {\sc Lamport,} Leslie.
%   \emph{\LaTeX: A Document Preparation System}.
%   2. vydání.
%   Massachusetts: Addison Wesley, 1994.
%   ISBN 0-201-52983-1.
%
% \end{thebibliography}


\listoffigures
\openright
\end{document}
